

Yes. The implication holds for every hyperfield. The key point is that a classical theorem of Krasner says that every hyperfield is a quotient of a field by a subgroup of its multiplicative group; then the statement lifts to the ordinary field identity
\[
X^3+Y^3+Z^3-3XYZ=(X+Y+Z)(X^2+Y^2+Z^2-XY-YZ-ZX).
\]
Relevant background on hyperfields and quotient/dequantization constructions includes
[ar5iv.labs.arxiv.org](https://ar5iv.labs.arxiv.org/html/2307.09400),
[ar5iv.labs.arxiv.org](https://ar5iv.labs.arxiv.org/html/1006.3034),
[arxiv.org](https://arxiv.org/html/2309.17302), and
[math.stonybrook.edu](https://www.math.stonybrook.edu/%7Eoleg/VersoTeX/Show/Hyperfields/hf.html).

Below is a compilable LaTeX document containing a complete proof, with the representation theorem cited carefully.

```latex
\documentclass{article}
\usepackage{amsmath}
\usepackage{amssymb}
\usepackage{amsthm}
\usepackage{hyperref}

\newtheorem{theorem}{Theorem}
\newtheorem{lemma}[theorem]{Lemma}
\newtheorem{definition}[theorem]{Definition}
\newtheorem{remark}[theorem]{Remark}

\title{A cubic zero-sum implication for hyperfields}
\author{}
\date{}

\begin{document}
\maketitle

\section*{Notation and conventions}

Let \(\mathbb H\) be a hyperfield with multivalued addition \(\boxplus\), additive identity \(0\), and multiplicative identity \(1\). For the definition of hyperfields and the extension of \(\boxplus\) to finite strings and subsets, see for instance
\href{https://ar5iv.labs.arxiv.org/html/2307.09400}{ar5iv.labs.arxiv.org}
and
\href{https://arxiv.org/html/2309.17302}{arxiv.org}.
For \(a\in\mathbb H\), \(-a\) denotes the unique additive inverse of \(a\), i.e.\ the unique element satisfying
\[
0\in a\boxplus(-a).
\]
In the expression \(-xyz\), the minus sign applies to the product: \(-xyz=-(xyz)\).

Repeated hyperaddition is understood associatively, as usual; equivalently one may flatten repeated hyperadditions using the power-set monad, as in
\href{https://ar5iv.labs.arxiv.org/html/2307.09400}{ar5iv.labs.arxiv.org}.

We shall prove the following theorem.

\begin{theorem}\label{thm:main}
Let \(\mathbb H\) be a hyperfield, and let \(x,y,z\in\mathbb H\). Suppose
\[
0\in x\boxplus y\boxplus z.
\]
Then
\[
0\in x^3\boxplus y^3\boxplus z^3\boxplus(-(xyz))\boxplus(-(xyz))\boxplus(-(xyz)).
\]
\end{theorem}

The proof uses Krasner's representation theorem for hyperfields and a lifting argument to an ordinary field.

\section{Quotient hyperfields}

\begin{definition}\label{def:quotient-hyperfield}
Let \(F\) be a field and let \(G\le F^\times\) be a subgroup of its multiplicative group. Write
\[
F/G:=\{0\}\cup F^\times/G.
\]
Multiplication is defined by
\[
0\cdot A=0,\qquad (aG)(bG)=abG.
\]
Hyperaddition is defined as follows. For nonzero cosets \(A=aG\) and \(B=bG\),
\[
A\boxplus B
=
\{\,cG\in F^\times/G : \exists \alpha\in A,\ \beta\in B\text{ with }\alpha+\beta=c\,\},
\]
and \(0\in A\boxplus B\) precisely when there exist \(\alpha\in A,\beta\in B\) with \(\alpha+\beta=0\). Also
\[
0\boxplus A=A\boxplus 0=\{A\},\qquad 0\boxplus 0=\{0\}.
\]
This is the usual quotient hyperfield construction; see
\href{https://ar5iv.labs.arxiv.org/html/1006.3034}{ar5iv.labs.arxiv.org}
and
\href{https://www.math.stonybrook.edu/\%7Eoleg/VersoTeX/Show/Hyperfields/hf.html}{math.stonybrook.edu}.
\end{definition}

The following lifting property is elementary but crucial.

\begin{lemma}\label{lem:lift-zero-sum}
Let \(F/G\) be a quotient hyperfield, and let \(A_1,\dots,A_n\in F/G\). Then
\[
0\in A_1\boxplus\cdots\boxplus A_n
\]
if and only if there exist representatives \(a_i\in A_i\) such that
\[
a_1+\cdots+a_n=0
\]
in the field \(F\). Here if \(A_i=0\), the only representative is \(a_i=0\).
\end{lemma}

\begin{proof}
For \(n=1\), the statement is clear. For \(n=2\), it is exactly the definition of hyperaddition in a quotient hyperfield.

Assume the result for \(n-1\). If
\[
0\in A_1\boxplus\cdots\boxplus A_n,
\]
then by associativity of hyperaddition there is some
\[
B\in A_1\boxplus\cdots\boxplus A_{n-1}
\]
such that \(0\in B\boxplus A_n\). By the case \(n=2\), there exist \(b\in B\) and \(a_n\in A_n\) with \(b+a_n=0\). By the induction hypothesis applied to \(B\), there exist \(a_i\in A_i\) for \(1\le i\le n-1\) such that
\[
a_1+\cdots+a_{n-1}=b.
\]
Thus
\[
a_1+\cdots+a_{n-1}+a_n=0.
\]

Conversely, suppose \(a_i\in A_i\) and \(a_1+\cdots+a_n=0\). Put \(b=a_1+\cdots+a_{n-1}\). By the induction hypothesis,
\[
b\in A_1\boxplus\cdots\boxplus A_{n-1}.
\]
Since \(b+a_n=0\), the definition of quotient hyperaddition gives
\[
0\in b\boxplus A_n.
\]
Hence
\[
0\in A_1\boxplus\cdots\boxplus A_n.
\]
\end{proof}

\section{Krasner's representation theorem}

We use the following classical theorem of Krasner. Expositions and discussions of this result can be found in
\href{https://ar5iv.labs.arxiv.org/html/1006.3034}{ar5iv.labs.arxiv.org},
\href{https://ar5iv.labs.arxiv.org/html/2307.09400}{ar5iv.labs.arxiv.org},
and
\href{https://www.math.stonybrook.edu/\%7Eoleg/VersoTeX/Show/Hyperfields/hf.html}{math.stonybrook.edu}.

\begin{theorem}[Krasner]\label{thm:krasner}
Every hyperfield \(\mathbb H\) is isomorphic to a quotient hyperfield \(F/G\), where \(F\) is a field and \(G\le F^\times\) is a subgroup of its multiplicative group.
\end{theorem}

Thus, without loss of generality, one may prove universal hyperfield identities by working in a quotient hyperfield \(F/G\), provided the argument uses only the hyperfield operations and their quotient description.

\section{The field identity}

\begin{lemma}\label{lem:field-identity}
Let \(F\) be a field, and let \(X,Y,Z\in F\) satisfy
\[
X+Y+Z=0.
\]
Then
\[
X^3+Y^3+Z^3+(-XYZ)+(-XYZ)+(-XYZ)=0.
\]
Equivalently,
\[
X^3+Y^3+Z^3-3XYZ=0.
\]
\end{lemma}

\begin{proof}
In the polynomial ring \(\mathbb Z[X,Y,Z]\), one has the universal identity
\[
X^3+Y^3+Z^3-3XYZ
=
(X+Y+Z)(X^2+Y^2+Z^2-XY-YZ-ZX).
\]
Reducing coefficients in the field \(F\), the same identity holds in \(F[X,Y,Z]\). If \(X+Y+Z=0\), the right-hand side is \(0\), hence
\[
X^3+Y^3+Z^3-3XYZ=0.
\]
But in the field \(F\),
\[
-3XYZ=(-XYZ)+(-XYZ)+(-XYZ),
\]
so the desired equality follows.
\end{proof}

\section{Proof of the main theorem}

\begin{proof}[Proof of Theorem~\ref{thm:main}]
By Krasner's representation theorem (Theorem~\ref{thm:krasner}), there exist a field \(F\), a subgroup \(G\le F^\times\), and a hyperfield isomorphism
\[
\phi:\mathbb H\longrightarrow F/G.
\]
Write
\[
A=\phi(x),\qquad B=\phi(y),\qquad C=\phi(z).
\]
Since \(\phi\) preserves hyperaddition and \(0\), the hypothesis gives
\[
0\in A\boxplus B\boxplus C.
\]
By Lemma~\ref{lem:lift-zero-sum}, there exist representatives
\[
X\in A,\qquad Y\in B,\qquad Z\in C
\]
such that
\[
X+Y+Z=0
\]
in \(F\).

By Lemma~\ref{lem:field-identity},
\[
X^3+Y^3+Z^3+(-XYZ)+(-XYZ)+(-XYZ)=0.
\]
Now observe that:
\begin{itemize}
\item \(X^3\in A^3=\phi(x^3)\);
\item \(Y^3\in B^3=\phi(y^3)\);
\item \(Z^3\in C^3=\phi(z^3)\);
\item \(-XYZ\in -(ABC)=\phi(-(xyz))\).
\end{itemize}
Indeed, multiplication in \(F/G\) is induced by multiplication in \(F\), and additive inverses in \(F/G\) are represented by negatives in \(F\).

Therefore, applying Lemma~\ref{lem:lift-zero-sum} to the six hyperfield elements
\[
A^3,\ B^3,\ C^3,\ -(ABC),\ -(ABC),\ -(ABC),
\]
we obtain
\[
0\in
A^3\boxplus B^3\boxplus C^3
\boxplus (-(ABC))\boxplus (-(ABC))\boxplus (-(ABC)).
\]
Finally, since \(\phi\) is a hyperfield isomorphism, pulling back by \(\phi^{-1}\) yields
\[
0\in
x^3\boxplus y^3\boxplus z^3
\boxplus (-(xyz))\boxplus (-(xyz))\boxplus (-(xyz)).
\]
This proves Theorem~\ref{thm:main}.
\end{proof}

\begin{remark}
The argument shows that the implication is not a special feature of any particular hyperfield; it is a consequence of Krasner's representation theorem together with the classical polynomial identity over fields. In the terminology sometimes used in the literature, this says that every hyperfield is \emph{perfect} with respect to this cubic relation.
\end{remark}

\begin{thebibliography}{9}

\bibitem{Factoring}
Factoring multivariate polynomials over hyperfields and the multivariable Descartes' problem,
\href{https://ar5iv.labs.arxiv.org/html/2307.09400}{ar5iv.labs.arxiv.org}.

\bibitem{Viro}
Hyperfields for Tropical Geometry I. Hyperfields and dequantization,
\href{https://ar5iv.labs.arxiv.org/html/1006.3034}{ar5iv.labs.arxiv.org}.

\bibitem{TropicalExtensions}
Geometry of tropical extensions of hyperfields,
\href{https://arxiv.org/html/2309.17302}{arxiv.org}.

\bibitem{StonyBrook}
Hyperfields and dequantization, lecture notes/web page,
\href{https://www.math.stonybrook.edu/\%7Eoleg/VersoTeX/Show/Hyperfields/hf.html}{math.stonybrook.edu}.

\end{thebibliography}

\end{document}
```